\cotefile{157-5}

\scanpage{573}

\mspage{482}

ce qui signifie que $x_\alpha$ n'appartient à aucun des $\deleted{\unsure{S_{\alpha}}}\underline S_\beta={}_{x_\beta}M$, pour $\beta\in J$. Donc $\underline S$ n'est pas réunion d'une petite sous-famille de la famille des $ {}_{x_\alpha}M$, donc en vertu de $\unsure{4.16.7.1|4.6.7.1}$, $\underline S$ n'est pas accessible dans $M$.

Inversement, soit $\underline S$ une telle sous-catégorie de $M$. Considérons un bon ordre sur $\mathrm{Ob}\,\underline S$ \inserted{(objets $s_\alpha$)}, et construisons la suite des $x_\alpha$ \inserted{$(\geq s_\alpha)$} par récurrence transfinie.

% slanted
\margin{\uline{NB} On a $\operatorname{card}\underline S\notin U$, \deleted{\unsure{sin}}sinon $\underline S$ serait petite, donc accessible dans $M$, contrairement à l'hypothèse.}

1) On choisit $\deleted{x_0\in\mathrm{Ob}\,\underline S}$ $x_0={}$ premier élément de $\mathrm{Ob}\,S$ (existe, car $\mathrm{Ob}\,\underline S\ne\varnothing$ i.e. $\underline S\ne\varnothing$, ou sinon $\underline S$ serait accessible dans $M$).

2) Si les $x_\beta$ ($\beta<\alpha$) sont construits, on prend pour $x_\alpha$ le premier élément de $\underline S$ \inserted{\unsure{majorant} les $x_\beta$ et} qui n'appartient pas à la sous-catégorie fermée de $M$ engendrée par les $x_\beta$ ($\beta<\alpha$). [Cette sous-catégorie fermée est $\subset\underline S$, et $\ne\underline S$ car $\underline S$ n'est pas accessible dans $M$] On a donc bien

\[
\operatorname{Hom}(x_\beta,x_\alpha)=\varnothing\quad\text{si }\beta<\alpha
\]

% archive sheet number: 572
